\section{Zaključek}
V prispevku smo obravnavali ravnotežje kolesno-nožnega robota na enem samem kolesu, ki predstavlja nestabilen, sklopljen in podaktuiran regulacijski problem z nagnjeno ravnovesno lego, ki ga s klasičnim modelom težko opišemo. Pokazali smo, da je takšen problem mogoče regulirati s strategijo učeno z globokim spodbujevanim učenjem, ki uporablja celice GRU kot osnovno arhitekturo. Strategija se je v mejnem območju učenja (začetni nagib $47^\circ \pm 11^\circ$) izkazala zanesljivo. Ker je robustno območje delovanja tako veliko, je tudi prehod iz dvokolesnega v enokolesno ravnovesje lažje. Rešitev namenoma deluje v okviru omejitev ciljne strojne opreme z majhno mrežo, fiksno normalizacijo, frekvenco regulacije $50$ Hz ter z randomizacijo domene, prilagojeno naključnosti realnega okolja, kar odpira pot neposrednemu prenosu na vgrajeni sistem v realnem okolju. 

Nadaljevanje tega dela definitivno odpira vrata za uporabo takšnih robotskih sistemov v realnih nestrukturiranih območjih, kjer bi lahko s pomočjo podobnih inherentno nestabilnih gibov omogočilo mobilnim robotom boljše premagovanje ovir. Uporaba celic GRU bi tako lahko bila uporabljena za boljše napovedovanje naslednjih stanj in premagovanju nestabilnostih v okolju samem.